\subsection{审查与修复表}
\begin{center}
\textbf{量化作业AI代码审查与修复表}
\end{center}

\subsubsection*{填写说明}
\begin{enumerate}
\item 本表为作业必填附件，需随策略代码、多轮交互日志一同提交；
\item 「代码出错位置」填写规范：策略文件路径 + 行号范围 + 核心错误代码片段，无错误填「无」；
\item 「严重等级」定义：
\begin{itemize}
\item 高（P0）：致命缺陷，直接导致策略逻辑失效/\allowbreak{}回测失真/\allowbreak{}实盘亏损，必须修复
\item 中（P1）：核心缺陷，影响策略稳定性/\allowbreak{}收益真实性，建议修复
\item 低（P2）：优化类问题，不影响核心逻辑，可选择性修复
\end{itemize}
\item 「修复情况」选项：未修复 /\allowbreak{} 部分修复 /\allowbreak{} 已修复，需同步填写「修复说明」。
\end{enumerate}
\subsubsection*{一、基础信息填写区}
\begingroup
\small
\setlength{\tabcolsep}{3pt}
\renewcommand{\arraystretch}{1.16}
\begin{longtable}{@{}L{0.24\textwidth}L{0.70\textwidth}@{}}
\toprule
\textbf{项目} & \textbf{填写内容} \\
\midrule
\endfirsthead
\toprule
\textbf{项目} & \textbf{填写内容} \\
\midrule
\endhead
\bottomrule
\endfoot
\bottomrule
\endlastfoot
学生姓名 & 丁致宇 \\
学号 & 202331060205 \\
作业名称 & 期末大作业：行业轮动多因子量化策略 \\
策略周期 & \checkedbox{} 混合周期（月度调仓，日频回测）  \uncheckedbox{} 日频短线  \uncheckedbox{} 周频中线 \\
代码来源 & \checkedbox{} AI生成+人工修改  \uncheckedbox{} 纯AI生成  \uncheckedbox{} 人工自主编写 \\
AI生成内容占比 & 70\%（AI辅助生成代码与报告，人工审查、口径确认和修复） \\
策略核心因子类型 & \checkedbox{} 快变量为主+慢变量过滤（动量/\allowbreak{}成交额为快变量，PB作估值过滤） \\
\end{longtable}
\endgroup
\subsubsection*{二、量化策略AI代码核心审查与修复表（主表）}
\subsubsection*{模块1：时序逻辑与未来函数专项审查（AI高频重灾区，P0级核心）}
\begingroup
\tiny
\setlength{\tabcolsep}{1pt}
\renewcommand{\arraystretch}{1.16}
\begin{longtable}{@{}L{0.045\textwidth}L{0.190\textwidth}L{0.140\textwidth}L{0.190\textwidth}L{0.060\textwidth}L{0.060\textwidth}L{0.205\textwidth}@{}}
\toprule
\textbf{序号} & \textbf{审查项} & \textbf{代码出错位置} & \textbf{问题详细描述} & \textbf{严重等级} & \textbf{修复情况} & \textbf{修复说明（含修复后代码片段）} \\
\midrule
\endfirsthead
\toprule
\textbf{序号} & \textbf{审查项} & \textbf{代码出错位置} & \textbf{问题详细描述} & \textbf{严重等级} & \textbf{修复情况} & \textbf{修复说明（含修复后代码片段）} \\
\midrule
\endhead
\bottomrule
\endfoot
\bottomrule
\endlastfoot
1.1 & 无未来函数：交易信号仅使用历史数据，禁止用当日/\allowbreak{}未来数据计算当日信号 & final\_\allowbreak{}project/\allowbreak{}factors.py:90-127；main.py:45,73 & 信号日因子仅使用月末及之前的价格、PB、成交额；样本按 fwd\_\allowbreak{}date 划分，避免未来收益进入训练。 & 高(P0) & 已修复 & 修复后保持 fwd\_\allowbreak{}return=close[d\_\allowbreak{}next]/\allowbreak{}close[d]-1，并用 filter\_\allowbreak{}panel(..., by="fwd\_\allowbreak{}date") 划分训练/\allowbreak{}回测。 \\
1.2 & 时间戳严格递增：数据加载、信号计算、交易执行按时间顺序执行，无乱序访问 & final\_\allowbreak{}project/\allowbreak{}data\_\allowbreak{}loader.py:82-92 & 三份CSV先转 Timestamp，再按共同交易日升序对齐，月末交易日按日历月最大交易日生成。 & 低(P2) & 已修复 & common\_\allowbreak{}dates 排序后统一索引，month\_\allowbreak{}ends=\_\allowbreak{}month\_\allowbreak{}end\_\allowbreak{}dates(common\_\allowbreak{}dates)。 \\
1.3 & 信号与交易分离：今日计算的信号，仅用于次日/\allowbreak{}下一个周期交易，无「当日信号当日成交」 & final\_\allowbreak{}project/\allowbreak{}backtest.py:94-125；config.py:88 & 配套申万行业数据无 open 字段，无法严格次日开盘成交；使用月末收盘到次月末收盘的代理口径。 & 中(P1) & 已修复 & 已在报告中明确“行业指数仅有收盘价，以月末收盘执行”；交易日收益只取 (d,d\_\allowbreak{}next]。 \\
1.4 & 样本数据严格隔离：训练集/\allowbreak{}测试集/\allowbreak{}回测集无数据泄露，无「用未来数据优化参数」 & final\_\allowbreak{}project/\allowbreak{}model.py:20-74；main.py:61-68 & 原 GroupKFold 仅防同月截面泄漏，不保证训练月早于验证月；质检保留因子原未接入模型。 & 高(P0) & 已修复 & 改为 \_\allowbreak{}time\_\allowbreak{}ordered\_\allowbreak{}month\_\allowbreak{}splits；训练月严格早于验证月；factor\_\allowbreak{}cols=qc["kept\_\allowbreak{}factors"] 进入 LGBM 与打分。 \\
1.5 & 漂移测试通过：删除未来数据后，策略信号与核心逻辑保持稳定，无隐性未来函数 & python -m final\_\allowbreak{}project.main --build-docx & 删除未来月份后，主流程可复现；2025样本外仍只使用2024-12及以前信号生成1月持仓。 & 中(P1) & 已修复 & 已重新运行完整流程，输出 metrics\_\allowbreak{}primary.csv 与 metrics\_\allowbreak{}final\_\allowbreak{}risk\_\allowbreak{}control.csv。 \\
\end{longtable}
\endgroup
\subsubsection*{模块2：因子与数据处理专项审查（高频快变量专属）}
\begingroup
\tiny
\setlength{\tabcolsep}{1pt}
\renewcommand{\arraystretch}{1.16}
\begin{longtable}{@{}L{0.045\textwidth}L{0.190\textwidth}L{0.140\textwidth}L{0.190\textwidth}L{0.060\textwidth}L{0.060\textwidth}L{0.205\textwidth}@{}}
\toprule
\textbf{序号} & \textbf{审查项} & \textbf{代码出错位置} & \textbf{问题详细描述} & \textbf{严重等级} & \textbf{修复情况} & \textbf{修复说明（含修复后代码片段）} \\
\midrule
\endfirsthead
\toprule
\textbf{序号} & \textbf{审查项} & \textbf{代码出错位置} & \textbf{问题详细描述} & \textbf{严重等级} & \textbf{修复情况} & \textbf{修复说明（含修复后代码片段）} \\
\midrule
\endhead
\bottomrule
\endfoot
\bottomrule
\endlastfoot
2.1 & 快变量筛选合理：仅使用与策略周期匹配的日/\allowbreak{}周频快变量，无周期错配 & final\_\allowbreak{}project/\allowbreak{}factors.py:36-56 & 使用日频收盘价、成交额和PB计算月度信号，周期与月度轮动匹配。 & 低(P2) & 已修复 & 20/\allowbreak{}60日动量、20日成交额占比均为日频快变量汇总到月末截面。 \\
2.2 & 慢变量使用合规：慢变量仅用于风险过滤/\allowbreak{}风格筛选，不直接作为短线交易主信号 & final\_\allowbreak{}project/\allowbreak{}factors.py:43-45；report.py & PB为估值变量，不直接做短线主信号，而是经Z-Score和IC方向参与综合打分。 & 低(P2) & 已修复 & PB方向由样本内IC确定；低PB得分更高。 \\
2.3 & 高频数据去噪完整：异常值剔除、滑动平均、滤波处理等，无「用单日极端数据做因子」 & final\_\allowbreak{}project/\allowbreak{}factors.py:105-112 & 早期回溯不足和缺失值已剔除；截面Z-Score消除量纲差异。 & 低(P2) & 已修复 & panel.dropna(...) 后按 signal\_\allowbreak{}date 分组做 \_\allowbreak{}zscore\_\allowbreak{}cross\_\allowbreak{}section。 \\
2.4 & 无因子共线性：相关系数\textgreater{}0.7的同源因子已剔除，无AI堆砌的重复逻辑因子 & final\_\allowbreak{}project/\allowbreak{}quality.py:72-98；main.py:61-68 & 相关性最高约0.53，低于阈值；原质检结果未进入训练路径的问题已修复。 & 中(P1) & 已修复 & kept\_\allowbreak{}factors 传入 train\_\allowbreak{}lgbm /\allowbreak{} ic\_\allowbreak{}directions /\allowbreak{} composite\_\allowbreak{}score，避免质检流于形式。 \\
2.5 & 因子经济逻辑清晰：所有因子均有明确的金融/\allowbreak{}交易逻辑支撑，无纯数学组合的无效因子 & final\_\allowbreak{}project/\allowbreak{}config.py:52-71；report.py & 四类因子均对应景气、估值、动量、资金关注度，金融含义清晰。 & 低(P2) & 已修复 & 报告2.1/\allowbreak{}2.2中写明因子口径、方向和经济解释。 \\
\end{longtable}
\endgroup
\subsubsection*{模块3：交易与风控合规专项审查（班级规则强制要求）}
\begingroup
\tiny
\setlength{\tabcolsep}{1pt}
\renewcommand{\arraystretch}{1.16}
\begin{longtable}{@{}L{0.045\textwidth}L{0.190\textwidth}L{0.140\textwidth}L{0.190\textwidth}L{0.060\textwidth}L{0.060\textwidth}L{0.205\textwidth}@{}}
\toprule
\textbf{序号} & \textbf{审查项} & \textbf{代码出错位置} & \textbf{问题详细描述} & \textbf{严重等级} & \textbf{修复情况} & \textbf{修复说明（含修复后代码片段）} \\
\midrule
\endfirsthead
\toprule
\textbf{序号} & \textbf{审查项} & \textbf{代码出错位置} & \textbf{问题详细描述} & \textbf{严重等级} & \textbf{修复情况} & \textbf{修复说明（含修复后代码片段）} \\
\midrule
\endhead
\bottomrule
\endfoot
\bottomrule
\endlastfoot
3.1 & 交易成本完整：已嵌入佣金+印花税+滑点（综合成本≥0.3\%），无零成本理想化回测 & final\_\allowbreak{}project/\allowbreak{}backtest.py:99-125,155-216 & 基线调仓成本已计入；原回撤减仓未计入额外换手成本，已补。 & 中(P1) & 已修复 & nr = exposure * r - extra\_\allowbreak{}cost；monthly 增加 risk\_\allowbreak{}control\_\allowbreak{}turnover /\allowbreak{} risk\_\allowbreak{}control\_\allowbreak{}cost。 \\
3.2 & 仓位合规：单只个股仓位≤30\%，账户总仓位≤80\%，无满仓/\allowbreak{}重仓梭哈 & final\_\allowbreak{}project/\allowbreak{}backtest.py:40-60；config.py:83-85 & 单行业上限30\%；Top5等权20\%不触发，Top3超额部分留现金。本题要求行业组合100\%配置，账户80\%规则不适用。 & 低(P2) & 已修复 & w=np.minimum(w, cap)，不向上重分配，超出部分保留现金。 \\
3.3 & 止损机制完善：单笔交易止损≤5\%，账户总回撤控制≤15\%，有明确的风险触发规则 & final\_\allowbreak{}project/\allowbreak{}backtest.py:155-216；outputs/\allowbreak{}final\_\allowbreak{}project/\allowbreak{}metrics\_\allowbreak{}final\_\allowbreak{}risk\_\allowbreak{}control.csv & 组合回撤触发减仓已实证；扣除新增成本后最终风控版年化14.58\%、最大回撤-9.87\%。 & 高(P0) & 已修复 & trig-5\%\_\allowbreak{}exp40\% 同时满足年化12\%-18\%和回撤≤10\%。 \\
3.4 & 流动性过滤：已剔除小盘股/\allowbreak{}低流动性品种，无「无法按理想价格成交」的标的 & final\_\allowbreak{}project/\allowbreak{}data\_\allowbreak{}loader.py:58-66；industry\_\allowbreak{}codes.py & 标的是申万一级行业指数，不直接交易小盘个股；流动性风险在报告中说明为ETF代理风险。 & 低(P2) & 已修复 & 31个行业代码固定映射，不做不可交易个股扫描。 \\
3.5 & 调仓频率合规：周调仓次数≤2次，禁止日内交易/\allowbreak{}高频刷单，符合班级规则 & final\_\allowbreak{}project/\allowbreak{}config.py:83-88；backtest.py:94-95 & 策略为月度调仓，不进行日内/\allowbreak{}高频交易；持有区间为下一月交易日。 & 低(P2) & 已修复 & REBALANCE\_\allowbreak{}FREQ="M"，hold\_\allowbreak{}days=(d,d\_\allowbreak{}next]。 \\
\end{longtable}
\endgroup
\subsubsection*{模块4：AI代码特有缺陷专项审查}
\begingroup
\tiny
\setlength{\tabcolsep}{1pt}
\renewcommand{\arraystretch}{1.16}
\begin{longtable}{@{}L{0.045\textwidth}L{0.190\textwidth}L{0.140\textwidth}L{0.190\textwidth}L{0.060\textwidth}L{0.060\textwidth}L{0.205\textwidth}@{}}
\toprule
\textbf{序号} & \textbf{审查项} & \textbf{代码出错位置} & \textbf{问题详细描述} & \textbf{严重等级} & \textbf{修复情况} & \textbf{修复说明（含修复后代码片段）} \\
\midrule
\endfirsthead
\toprule
\textbf{序号} & \textbf{审查项} & \textbf{代码出错位置} & \textbf{问题详细描述} & \textbf{严重等级} & \textbf{修复情况} & \textbf{修复说明（含修复后代码片段）} \\
\midrule
\endhead
\bottomrule
\endfoot
\bottomrule
\endlastfoot
4.1 & 无AI幻觉：无虚构的API/\allowbreak{}函数/\allowbreak{}库调用，所有import的依赖真实存在且版本兼容 & requirements.txt；python -m compileall -q final\_\allowbreak{}project & 未发现虚构API；依赖 pandas/\allowbreak{}numpy/\allowbreak{}scipy/\allowbreak{}sklearn/\allowbreak{}lightgbm/\allowbreak{}python-docx 均可导入。 & 低(P2) & 已修复 & 已通过 compileall 和完整主流程运行。 \\
4.2 & 无逻辑矛盾：条件判断无冲突，交易方向明确，无「同时买入和卖出」的悖论代码 & final\_\allowbreak{}project/\allowbreak{}main.py:61-68；robustness.py:73-81 & 原质检与模型路径存在断点；方法敏感性已显式使用同一 factor\_\allowbreak{}cols。 & 中(P1) & 已修复 & composite\_\allowbreak{}score(..., model\_\allowbreak{}result["factor\_\allowbreak{}cols"])。 \\
4.3 & 无过度拟合：策略在样本外数据表现稳定，无「回测收益极高、实盘持续亏损」的过拟合问题 & final\_\allowbreak{}project/\allowbreak{}model.py:20-74；outputs/\allowbreak{}final\_\allowbreak{}project/\allowbreak{}robustness\_\allowbreak{}time.csv & 原CV可能时序泄漏；策略2024跑输、2025跑赢，报告如实呈现机制依赖。 & 高(P0) & 已修复 & 扩展窗口CV + 2024/\allowbreak{}2025时间敏感性测试，未掩盖样本外风险。 \\
4.4 & 无硬编码参数：关键参数（如持仓周期、止损比例、因子窗口）可配置，无AI固定死的魔法数字 & final\_\allowbreak{}project/\allowbreak{}config.py & 关键窗口、阈值、成本、TopN、风控目标均集中配置，无散落魔法数字。 & 低(P2) & 已修复 & MOM\_\allowbreak{}MID\_\allowbreak{}WINDOW、COST\_\allowbreak{}RATE、TOP\_\allowbreak{}N\_\allowbreak{}DEFAULT 等统一管理。 \\
4.5 & 无冗余代码：AI生成的重复逻辑、无效循环、未使用的变量已删除，代码简洁高效 & final\_\allowbreak{}project/\allowbreak{}report.py:338-365；build\_\allowbreak{}docx & 原AI审核/\allowbreak{}交互记录字符串截断，Word报告未内嵌AI记录；已修复。 & 中(P1) & 已修复 & 新增 \_\allowbreak{}ai\_\allowbreak{}review\_\allowbreak{}rows /\allowbreak{} \_\allowbreak{}ai\_\allowbreak{}interaction\_\allowbreak{}items，Markdown与Word共用同一数据源。 \\
\end{longtable}
\endgroup
\subsubsection*{模块5：代码可读性与可维护性审查}
\begingroup
\tiny
\setlength{\tabcolsep}{1pt}
\renewcommand{\arraystretch}{1.16}
\begin{longtable}{@{}L{0.045\textwidth}L{0.190\textwidth}L{0.140\textwidth}L{0.190\textwidth}L{0.060\textwidth}L{0.060\textwidth}L{0.205\textwidth}@{}}
\toprule
\textbf{序号} & \textbf{审查项} & \textbf{代码出错位置} & \textbf{问题详细描述} & \textbf{严重等级} & \textbf{修复情况} & \textbf{修复说明（含修复后代码片段）} \\
\midrule
\endfirsthead
\toprule
\textbf{序号} & \textbf{审查项} & \textbf{代码出错位置} & \textbf{问题详细描述} & \textbf{严重等级} & \textbf{修复情况} & \textbf{修复说明（含修复后代码片段）} \\
\midrule
\endhead
\bottomrule
\endfoot
\bottomrule
\endlastfoot
5.1 & 注释完整：核心逻辑、因子计算、交易规则、风控机制均有清晰的中文注释 & final\_\allowbreak{}project/\allowbreak{}*.py & 核心函数均有中文docstring和关键注释；修复后补充风控成本说明。 & 低(P2) & 已修复 & apply\_\allowbreak{}drawdown\_\allowbreak{}control 注明额外换手按 cost\_\allowbreak{}rate 扣费。 \\
5.2 & 命名规范：变量、函数、类名符合行业标准，见名知意，无AI生成的无意义命名 & final\_\allowbreak{}project/\allowbreak{}*.py & 变量名围绕 market/\allowbreak{}panel/\allowbreak{}factor/\allowbreak{}score/\allowbreak{}nav/\allowbreak{}metrics，含义清晰。 & 低(P2) & 已修复 & 未发现无意义命名。 \\
5.3 & 结构清晰：代码按「数据加载→因子计算→信号生成→交易执行→风控→回测」模块化拆分，无单文件堆砌 & final\_\allowbreak{}project/\allowbreak{} & 已按数据加载、因子、质检、模型、回测、鲁棒性、报告拆分模块。 & 低(P2) & 已修复 & main.py 负责流程编排。 \\
5.4 & 错误处理完善：有完整的try-catch异常捕获，网络/\allowbreak{}数据/\allowbreak{}计算异常有明确的处理逻辑 & final\_\allowbreak{}project/\allowbreak{}data\_\allowbreak{}loader.py:82-88；factors.py:81-88 & 对齐和因子取值有防御性处理；无网络下载依赖。 & 低(P2) & 已修复 & common\_\allowbreak{}dates 交集对齐；KeyError月份跳过。 \\
5.5 & 结果可复现：随机种子固定，回测结果可100\%复现，无AI生成的随机波动代码 & final\_\allowbreak{}project/\allowbreak{}config.py:116-130；python -m final\_\allowbreak{}project.main --build-docx & 随机种子固定，输出可复现；已重新生成CSV、图表、Markdown与Word。 & 低(P2) & 已修复 & LGBM random\_\allowbreak{}state=42；完整流程运行成功。 \\
\end{longtable}
\endgroup
\clearpage
\subsubsection*{三、AI代码整体评估与总结区}
\begingroup
\scriptsize
\setlength{\tabcolsep}{2pt}
\renewcommand{\arraystretch}{1.16}
\begin{longtable}{@{}L{0.28\textwidth}L{0.66\textwidth}@{}}
\toprule
\textbf{项目} & \textbf{填写内容} \\
\midrule
\endfirsthead
\toprule
\textbf{项目} & \textbf{填写内容} \\
\midrule
\endhead
\bottomrule
\endfoot
\bottomrule
\endlastfoot
本次AI生成代码的核心缺陷TOP3 & 1. GroupKFold不满足时序验证，可能造成验证泄漏。\newline{} 2. 因子质检结果未接入LGBM训练与打分。\newline{} 3. 风控减仓未计入额外交易成本，且AI审核/\allowbreak{}交互材料未写入Word报告。 \\
人工修复与优化的核心亮点 & 1. 改为按月份扩展窗口验证，训练月严格早于验证月。\newline{} 2. 使用 qc["kept\_\allowbreak{}factors"] 驱动建模，并输出最终风控版净值/\allowbreak{}指标。\newline{} 3. 风控成本、AI附录、老师审查表全部补齐，重跑生成Word报告。 \\
本次策略最终合规性结论 & \checkedbox{} 完全合规（所有P0/\allowbreak{}P1问题已修复；行业数据无open字段，已明确使用收盘价代理执行口径）\newline{} \uncheckedbox{} 基本合规（仅剩余低优先级优化问题）\newline{} \uncheckedbox{} 不合规（存在未修复的致命缺陷） \\
对AI在量化高频策略中局限性的总结反思 & AI在量化策略中容易把“能跑通”误认为“无泄漏、可交易”：典型风险包括时序切分不严格、质检结果未真正进入模型、风控成本理想化、报告与代码输出不一致。因此必须由人工逐项核查时间线、成本、样本隔离和提交材料完整性。 \\
\end{longtable}
\endgroup
\clearpage
